\documentclass[11pt,a4paper]{article}
\usepackage[utf8]{inputenc}
\usepackage[spanish,es-tabla]{babel}
\usepackage{amsmath,amsfonts,amssymb}
\usepackage{geometry}
\usepackage{graphicx}
\usepackage{xcolor}

\usepackage{enumitem}
\usepackage{tikz}

\geometry{margin=2.5cm}

\title{\textbf{Soluci\'on Examen 2023-2024\\TAF STAR - UEC B\\17/10/2023}}
\author{}
\date{}

\begin{document}

\maketitle

\section*{Transmisi\'on ASK sobre canal AWGN}

\textbf{Datos del problema:}
\begin{itemize}
    \item Elementos binarios independientes e id\'enticamente distribuidos
    \item $\Pr(m_i = 0) = \Pr(m_i = 1)$
    \item D\'ebito binario: $D_b = 1.6 \times 10^6$ bits/s
    \item Modulaci\'on: 4-ASK
    \item Frecuencia portadora: $\nu_0 = 2 \times 10^9$ Hz
\end{itemize}

\textbf{Se\~nal modulada:}
\[
x(t) = \sum_{k=0}^{L-1} a_k h(t-kT)\cos(2\pi\nu_0 t + \theta_0)
\]

con la condici\'on $\int_{\mathbb{R}} |h(t)|^2\, dt = 1$.

\textbf{Funci\'on de transferencia del filtro:}
\[
H(\nu) = \int_{\mathbb{R}} h(t)e^{-i2\pi\nu t}\, dt
\]

con $\begin{cases}
H(\nu) \neq 0, & \text{si } |\nu| < 6 \times 10^5 \\
H(\nu) = 0, & \text{en otro caso}
\end{cases}$

\textbf{Canal:} AWGN

\textbf{Se\~nal recibida:} $r(t) = x(t) + w(t)$

\textbf{Receptor:}

\begin{center}
\begin{tikzpicture}[node distance=2cm]
\node (input) {$r(t)$};
\node[draw, right of=input] (mult) {$\times \cos(2\pi\nu_0 t + \theta_0)$};
\node[draw, right of=mult, node distance=3cm] (filter) {$g(t)$};
\node[draw, right of=filter] (sampler) {$t_n$};
\node[draw, right of=sampler] (decision) {Decisi\'on};
\node[right of=decision] (output) {$\hat{a}_n$};

\draw[->] (input) -- (mult);
\draw[->] (mult) -- (filter);
\draw[->] (filter) -- (sampler);
\draw[->] (sampler) -- node[above] {$z_n$} (decision);
\draw[->] (decision) -- (output);
\end{tikzpicture}
\end{center}

\subsection*{a) Definir el alfabeto 4-ASK}

\textbf{Enunciado:} Definir el alfabeto 4-ASK.

\vspace{0.3cm}
\textbf{Soluci\'on:}

El alfabeto 4-ASK est\'a formado por 4 niveles de amplitud equiespaciados y sim\'etricos respecto al origen:

\[
\boxed{\mathcal{D} = \{-3, -1, +1, +3\}}
\]

\textbf{Propiedades:}
\begin{itemize}
    \item $M = 4$ s\'imbolos $\Rightarrow$ $\log_2(4) = 2$ bits por s\'imbolo
    \item Separaci\'on entre s\'imbolos adyacentes: $\Delta = 2$
    \item Amplitudes sim\'etricas: media cero
    \item Energ\'ia media: $\mathbb{E}[a_k^2] = \frac{1}{4}(9 + 1 + 1 + 9) = 5$
\end{itemize}



\subsection*{b) Calcular T}

\textbf{Enunciado:} Calcular $T$.

\vspace{0.3cm}
\textbf{Soluci\'on:}

La relaci\'on entre d\'ebito binario, rapidez de s\'imbolo y periodo s\'imbolo es:

\[
D_s = \frac{D_b}{\log_2(M)} \quad \text{y} \quad T = \frac{1}{D_s}
\]

Para 4-ASK: $M = 4$, por lo tanto:

\[
D_s = \frac{1.6 \times 10^6}{2} = 0.8 \times 10^6 \text{ s\'imbolos/s} = 800 \text{ kbauds}
\]

El periodo s\'imbolo es:

\[
\boxed{T = \frac{1}{D_s} = \frac{1}{0.8 \times 10^6} = 1.25 \times 10^{-6} \text{ s} = 1.25 \text{ μs}}
\]

\textbf{Interpretaci\'on:}
\begin{itemize}
    \item Cada s\'imbolo transporta 2 bits
    \item Se transmiten 800,000 s\'imbolos por segundo
    \item Cada s\'imbolo dura 1.25 microsegundos
\end{itemize}



\subsection*{c) Expresi\'on de la densidad espectral de potencia}

\textbf{Enunciado:} Dar la expresi\'on de la densidad espectral de potencia de $x(t)$.

\vspace{0.3cm}
\textbf{Soluci\'on:}

Para una modulaci\'on lineal sobre portadora con s\'imbolos independientes de media cero, la densidad espectral de potencia es:

\[
\boxed{S_x(\nu) = \frac{\sigma_a^2}{4T}\left[|H(\nu - \nu_0)|^2 + |H(\nu + \nu_0)|^2\right]}
\]

Donde:
\begin{itemize}
    \item $\sigma_a^2 = \mathbb{E}[a_k^2] = 5$ (varianza de los s\'imbolos, calculada anteriormente)
    \item $T = 1.25 \times 10^{-6}$ s (periodo s\'imbolo)
    \item $H(\nu)$ es la funci\'on de transferencia del filtro de conformaci\'on
    \item $\nu_0 = 2 \times 10^9$ Hz (frecuencia portadora)
\end{itemize}

Sustituyendo:

\[
\boxed{S_x(\nu) = \frac{5}{4 \times 1.25 \times 10^{-6}}\left[|H(\nu - \nu_0)|^2 + |H(\nu + \nu_0)|^2\right] = 10^6\left[|H(\nu - \nu_0)|^2 + |H(\nu + \nu_0)|^2\right]}
\]

\textbf{Interpretaci\'on:}
\begin{itemize}
    \item El espectro est\'a centrado en $\pm\nu_0$ (bandas laterales sim\'etricas)
    \item La forma del espectro est\'a determinada por $|H(\nu)|^2$
    \item El factor $\frac{1}{4T}$ proviene de la modulaci\'on sobre portadora con $\cos^2$ promedio igual a $\frac{1}{2}$
    \item La potencia total es proporcional a $\sigma_a^2$
\end{itemize}



\subsection*{d) Calcular la eficiencia espectral}

\textbf{Enunciado:} Calcular la eficiencia espectral.

\vspace{0.3cm}
\textbf{Soluci\'on:}

\textbf{Paso 1: Determinar la banda ocupada}

Del enunciado, sabemos que $H(\nu) \neq 0$ solo para $|\nu| < 6 \times 10^5$ Hz.

Esto significa que el filtro paso banda tiene banda base $[-6 \times 10^5, 6 \times 10^5]$ Hz.

Para la se\~nal modulada, la banda de frecuencias ocupada es:

\[
[\nu_0 - 6 \times 10^5, \nu_0 + 6 \times 10^5]
\]

Por lo tanto, la banda ocupada es:

\[
B_{occ} = 2 \times 6 \times 10^5 = 1.2 \times 10^6 \text{ Hz} = 1.2 \text{ MHz}
\]

\textbf{Paso 2: Calcular la eficiencia espectral}

La eficiencia espectral se define como:

\[
\eta = \frac{D_b}{B_{occ}} \text{ [bits/s/Hz]}
\]

Sustituyendo:

\[
\eta = \frac{1.6 \times 10^6}{1.2 \times 10^6} = \frac{1.6}{1.2} = \boxed{1.33 \text{ bits/s/Hz}}
\]

\textbf{An\'alisis:}

\begin{itemize}
    \item \textbf{Comparaci\'on con el l\'imite de Nyquist:} 
    
    La banda m\'inima te\'orica para transmisi\'on sin ISI es $B_{min} = \frac{D_s}{2} = \frac{0.8 \times 10^6}{2} = 0.4$ MHz (para modulaci\'on sobre portadora).
    
    Ratio: $\frac{B_{occ}}{B_{min}} = \frac{1.2}{0.4} = 3$
    
    El sistema utiliza 3 veces la banda m\'inima te\'orica.
    
    \item \textbf{Roll-off impl\'icito:} 
    
    Como $B_{occ} = D_s(1 + \beta)$ para modulaci\'on paso banda, tenemos:
    \[
    1.2 \times 10^6 = 0.8 \times 10^6 (1 + \beta) \Rightarrow \beta = 0.5
    \]
    
    \item \textbf{Eficiencia moderada:} El valor de 1.33 bits/s/Hz es moderado. Sistemas m\'as eficientes pueden alcanzar 4-6 bits/s/Hz con modulaciones de mayor orden.
\end{itemize}



\subsection*{e) Definir $g(t)$, $t_n$ y condici\'on TSWOC}

\textbf{Enunciado:} Definir $g(t)$, $t_n$ y la condici\'on sobre $h(t)$ tal que el TSWOC se satisfaga. Las expresiones de $g(t)$ y $t_n$ deben escribirse en funci\'on de los par\'ametros de la se\~nal modulada.

\vspace{0.3cm}
\textbf{Soluci\'on:}

\textbf{Filtro de recepci\'on $g(t)$:}

Para un receptor \'optimo de m\'axima verosimilitud, el filtro de recepci\'on debe ser el \textbf{filtro adaptado} al filtro de transmisi\'on:

\[
\boxed{g(t) = h(t_0 - t)}
\]

Donde $t_0$ es un retardo que garantiza la causalidad del filtro (depende del soporte de $h(t)$).

\textbf{Instantes de muestreo $t_n$:}

Los instantes \'optimos de muestreo son:

\[
\boxed{t_n = t_0 + nT}
\]

Donde:
\begin{itemize}
    \item $n \in \mathbb{Z}$ es el \'indice del s\'imbolo
    \item $T$ es el periodo s\'imbolo
    \item $t_0$ es el retardo del filtro adaptado
\end{itemize}

\textbf{Condici\'on TSWOC sobre $h(t)$:}

El criterio Time-Shifted-Waveform Orthogonality Criterion (TSWOC) requiere:

\[
\boxed{\int_{\mathbb{R}} h(t) h(t - kT)\, dt = \delta_{k,0} \int_{\mathbb{R}} |h(t)|^2\, dt}
\]

Como el enunciado especifica que $\int_{\mathbb{R}} |h(t)|^2\, dt = 1$, la condici\'on se simplifica a:

\[
\boxed{\int_{\mathbb{R}} h(t) h(t - kT)\, dt = \delta_{k,0}}
\]

\textbf{Formulaci\'on alternativa en frecuencia:}

En el dominio de Fourier, el TSWOC es equivalente a:

\[
\boxed{\frac{1}{T}\sum_p |H(\nu - p/T)|^2 = 1, \quad \forall \nu \in \mathbb{R}}
\]

O con la normalizaci\'on del enunciado:

\[
\boxed{\sum_p |H(\nu - p/T)|^2 = T, \quad \forall \nu \in \mathbb{R}}
\]

\textbf{Significado de cada expresi\'on:}

\begin{itemize}
    \item $g(t) = h(t_0 - t)$: El filtro de recepci\'on es la versi\'on retardada y ''volteada en el tiempo'' del filtro de transmisi\'on. Esto maximiza el SNR.
    
    \item $t_n = t_0 + nT$: Los muestras se toman en instantes equiespaciados con periodo $T$, despu\'es del retardo $t_0$ del filtro.
    
    \item TSWOC: Garantiza que las versiones desplazadas del pulso sean ortogonales, eliminando la interferencia entre s\'imbolos.
\end{itemize}



\subsection*{f) Consecuencias de la satisfacci\'on del TSWOC}

\textbf{Enunciado:} ¿Cu\'ales son las consecuencias de la satisfacci\'on del TSWOC para la detecci\'on?

\vspace{0.3cm}
\textbf{Soluci\'on:}

La satisfacci\'on del criterio TSWOC tiene tres consecuencias fundamentales para la detecci\'on:

\textbf{1. Detecci\'on s\'imbolo por s\'imbolo \'optima:}

La decisi\'on sobre cada s\'imbolo $a_k$ puede tomarse de forma \textbf{independiente}, sin necesidad de conocer los s\'imbolos adyacentes $\{a_n\}_{n \neq k}$.

\textit{Justificaci\'on:} La muestra $z_k$ solo depende de $a_k$ y del ruido, no de otros s\'imbolos.

\textbf{2. Ausencia de interferencia entre s\'imbolos (ISI):}

En los instantes de muestreo $t_n = t_0 + nT$, no hay contribuci\'on de s\'imbolos vecinos.

\textit{Matem\'aticamente:} Si definimos $q(t) = (h * g)(t)$, entonces:
\[
q(t_0 + kT) = \int_{\mathbb{R}} h(\tau) g(t_0 + kT - \tau)\, d\tau = \delta_{k,0}
\]

Esto significa:
\begin{itemize}
    \item $q(t_0) = 1$ (el s\'imbolo $k=0$ contribuye completamente)
    \item $q(t_0 + kT) = 0$ para $k \neq 0$ (los dem\'as s\'imbolos no contribuyen)
\end{itemize}

\textbf{3. M\'aximo SNR a la salida del filtro:}

El filtro adaptado $g(t) = h(t_0 - t)$ maximiza la relaci\'on se\~nal a ruido en los instantes de muestreo.

\textit{SNR \'optimo:}
\[
\text{SNR} = \frac{|a_k|^2 \int_{\mathbb{R}} |h(t)|^2\, dt}{N_0/2} = \frac{|a_k|^2}{N_0/2}
\]

(usando la normalizaci\'on $\int |h(t)|^2\, dt = 1$)

\textbf{Resumen:}

\begin{center}
\begin{tabular}{|l|l|}
\hline
\textbf{Consecuencia} & \textbf{Beneficio} \\
\hline
Detecci\'on independiente & Simplifica el receptor \\
Ausencia de ISI & Evita propagaci\'on de errores \\
SNR m\'aximo & Minimiza probabilidad de error \\
\hline
\end{tabular}
\end{center}

\textbf{Conclusi\'on pr\'actica:} El receptor puede implementarse como un simple comparador de umbrales despu\'es del filtrado y muestreo, sin necesidad de ecualizadores o detectores de secuencia complejos.



\subsection*{g) Expresi\'on de $z_n$ asumiendo TSWOC satisfecho}

\textbf{Enunciado:} Asumiendo que el TSWOC est\'a satisfecho, calcular la expresi\'on de la entrada del dispositivo de decisi\'on $z_n$ y definir las reglas de decisi\'on sobre los s\'imbolos.

\vspace{0.3cm}
\textbf{Soluci\'on:}

\textbf{C\'alculo de $z_n$:}

\textbf{Paso 1:} Demodulaci\'on

\begin{align*}
y(t) &= r(t) \cos(2\pi\nu_0 t + \theta_0) \\
&= \left[\sum_{k=0}^{L-1} a_k h(t-kT)\cos(2\pi\nu_0 t + \theta_0) + w(t)\right] \cos(2\pi\nu_0 t + \theta_0)
\end{align*}

Usando $\cos^2(\alpha) = \frac{1 + \cos(2\alpha)}{2}$:

\[
y(t) = \frac{1}{2}\sum_{k=0}^{L-1} a_k h(t-kT) + \frac{1}{2}\sum_{k=0}^{L-1} a_k h(t-kT)\cos(4\pi\nu_0 t + 2\theta_0) + w_c(t)
\]

\textbf{Paso 2:} Filtrado

La componente de alta frecuencia centrada en $2\nu_0$ es eliminada por el filtro paso-bajo $g(t)$.

La salida del filtro es:
\[
z(t) = (y * g)(t) = \frac{1}{2}\sum_{k=0}^{L-1} a_k (h * g)(t-kT) + \gamma(t)
\]

donde $\gamma(t) = (w_c * g)(t)$ es el ruido filtrado.

\textbf{Paso 3:} Muestreo en $t_n = t_0 + nT$

\begin{align*}
z_n &= z(t_0 + nT) \\
&= \frac{1}{2}\sum_{k=0}^{L-1} a_k \int_{\mathbb{R}} h(\tau) g(t_0 + nT - kT - \tau)\, d\tau + \gamma_n \\
&= \frac{1}{2}\sum_{k=0}^{L-1} a_k \int_{\mathbb{R}} h(\tau) h(t_0 - (t_0 + nT - kT) + \tau)\, d\tau + \gamma_n \\
&= \frac{1}{2}\sum_{k=0}^{L-1} a_k \int_{\mathbb{R}} h(\tau) h(\tau + (k-n)T)\, d\tau + \gamma_n
\end{align*}

\textbf{Paso 4:} Aplicaci\'on del TSWOC

Como $\int_{\mathbb{R}} h(\tau) h(\tau + mT)\, d\tau = \delta_{m,0}$, solo el t\'ermino con $k = n$ sobrevive:

\[
\boxed{z_n = \frac{1}{2} a_n \int_{\mathbb{R}} |h(\tau)|^2\, d\tau + \gamma_n = \frac{1}{2} a_n + \gamma_n}
\]

(usando $\int |h(t)|^2\, dt = 1$)

Donde:
\begin{itemize}
    \item $a_n \in \{-3, -1, +1, +3\}$ es el s\'imbolo transmitido
    \item $\gamma_n$ es ruido gaussiano con media cero y varianza $\sigma^2 = \frac{N_0}{2}$
\end{itemize}

\textbf{Reglas de decisi\'on sobre s\'imbolos:}

Definiendo $\beta = \frac{1}{2}$, la constelaci\'on recibida (sin ruido) es:
\[
\left\{-\frac{3}{2}, -\frac{1}{2}, +\frac{1}{2}, +\frac{3}{2}\right\}
\]

Las fronteras de decisi\'on est\'an en $-1, 0, +1$.

\textbf{Regla de decisi\'on ML:}
\[
\boxed{
\hat{a}_n = \begin{cases}
-3 & \text{si } z_n < -1 \\
-1 & \text{si } -1 \leq z_n < 0 \\
+1 & \text{si } 0 \leq z_n < 1 \\
+3 & \text{si } z_n \geq 1
\end{cases}
}
\]

\textbf{Expresi\'on compacta:}
\[
\hat{a}_n = 2 \left\lfloor z_n + \frac{1}{2} \right\rfloor + 1 \quad \text{(para valores impares)}
\]

O mediante cuantificaci\'on de 4 niveles con fronteras en $\{-1, 0, +1\}$.



\newpage
\section*{Preguntas con respuestas cortas}

\subsection*{a) Principio del Gray mapping}

\textbf{Enunciado:} Sea la modulaci\'on 16-PSK cuyo alfabeto se muestra en la figura. Explicar el principio del Gray mapping.

\vspace{0.3cm}
\textbf{Soluci\'on:}

El \textbf{principio del Gray mapping} es una t\'ecnica de asignaci\'on de c\'odigos binarios a s\'imbolos de modulaci\'on que minimiza el impacto de los errores de s\'imbolo en la tasa de error binaria.

\textbf{Definici\'on formal:}

Un mapping binario a s\'imbolo satisface el criterio de Gray si:

\[
\boxed{\text{Dos s\'imbolos adyacentes en la constelaci\'on difieren en exactamente un bit}}
\]

\textbf{Caracter\'isticas del Gray mapping:}

\begin{enumerate}
    \item \textbf{Distancia de Hamming unitaria entre vecinos:}
    
    Si $s_i$ y $s_j$ son s\'imbolos adyacentes con c\'odigos binarios $c_i$ y $c_j$, entonces la distancia de Hamming:
    \[
    d_H(c_i, c_j) = 1
    \]
    
    \item \textbf{Minimizaci\'on de errores binarios:}
    
    En canal AWGN, los errores m\'as probables son confusiones entre s\'imbolos adyacentes (m\'as cercanos). Con Gray mapping, cada error de s\'imbolo causa t\'ipicamente solo un error de bit.
    
    \item \textbf{Construcci\'on:}
    
    Para modulaciones M-arias, los c\'odigos Gray se construyen de forma que al recorrer la constelaci\'on en orden, cada transici\'on cambie solo un bit.
\end{enumerate}

\textbf{Ejemplo para 16-PSK:}

Partiendo de un s\'imbolo con c\'odigo 0000, los s\'imbolos adyacentes podr\'ian ser:
\begin{itemize}
    \item Sentido horario: 0001 (cambio en bit 3)
    \item Sentido antihorario: 1000 (cambio en bit 0)
\end{itemize}

\textbf{Ventaja cuantitativa:}

Para SNR moderada a alta, la probabilidad de error binaria con Gray mapping es aproximadamente:
\[
P_b \approx \frac{1}{\log_2(M)} P_s
\]

Sin Gray mapping, podr\'ia ser hasta $\frac{\log_2(M)}{2} \cdot \frac{1}{\log_2(M)} P_s = \frac{1}{2} P_s$ en el peor caso.

\textbf{Conclusi\'on:} El Gray mapping es una estrategia simple pero eficaz para mejorar el rendimiento de sistemas de comunicaci\'on digital, especialmente en reg\'imenes de SNR medio-alto.



\subsection*{a-ii) Propuesta de Gray mapping para 16-PSK}

\textbf{Enunciado:} Proponer y escribir directamente en la figura un mapping binario a s\'imbolo que satisfaga el criterio de Gray mapping (es decir, asociar un c\'odigo binario para cada s\'imbolo).

\vspace{0.3cm}
\textbf{Soluci\'on:}

Para 16-PSK, los 16 s\'imbolos est\'an distribuidos uniformemente en un c\'irculo con \'angulos:
\[
\theta_k = \frac{2\pi k}{16} = \frac{\pi k}{8}, \quad k = 0, 1, 2, \ldots, 15
\]

\textbf{Gray mapping propuesto:}

\begin{center}
\begin{tabular}{|c|c|c|}
\hline
\textbf{$k$} & \textbf{\'Angulo $\theta_k$} & \textbf{C\'odigo Gray} \\
\hline
0 & $0°$ & 0000 \\
1 & $22.5°$ & 0001 \\
2 & $45°$ & 0011 \\
3 & $67.5°$ & 0010 \\
4 & $90°$ & 0110 \\
5 & $112.5°$ & 0111 \\
6 & $135°$ & 0101 \\
7 & $157.5°$ & 0100 \\
8 & $180°$ & 1100 \\
9 & $202.5°$ & 1101 \\
10 & $225°$ & 1111 \\
11 & $247.5°$ & 1110 \\
12 & $270°$ & 1010 \\
13 & $292.5°$ & 1011 \\
14 & $315°$ & 1001 \\
15 & $337.5°$ & 1000 \\
\hline
\end{tabular}
\end{center}

\textbf{Verificaci\'on del criterio Gray:}

\begin{itemize}
    \item $0000 \rightarrow 0001$: cambia bit 3 ✓
    \item $0001 \rightarrow 0011$: cambia bit 2 ✓
    \item $0011 \rightarrow 0010$: cambia bit 3 ✓
    \item $0010 \rightarrow 0110$: cambia bit 1 ✓
    \item Y as\'i sucesivamente alrededor del c\'irculo...
    \item $1000 \rightarrow 0000$ (cierre circular): cambia bit 0 ✓
\end{itemize}

\textbf{Patr\'on observado:}

El mapping sigue un c\'odigo Gray circular de 4 bits, donde:
\begin{itemize}
    \item Los dos bits m\'as significativos codifican el cuadrante (I, II, III, IV)
    \item Los dos bits menos significativos codifican la posici\'on dentro del cuadrante
    \item Cada transici\'on angular de $22.5°$ cambia exactamente un bit
\end{itemize}



\subsection*{b) Fronteras de decisi\'on y regiones para 16-APSK}

\textbf{Enunciado:} La constelaci\'on 16-APSK recibida sin ruido se muestra en la figura.

\textbf{b-i)} ¿C\'omo se define la frontera de decisi\'on entre dos s\'imbolos adyacentes para asegurar recepci\'on ML sobre un canal AWGN?

\textbf{b-ii)} Dibujar en la figura las regiones de decisi\'on a partir de las fronteras de decisi\'on que satisfacen esta regla.

\vspace{0.3cm}
\textbf{Soluci\'on:}

\textbf{b-i) Definici\'on de la frontera de decisi\'on:}

Para un receptor de m\'axima verosimilitud (ML) sobre un canal AWGN, la frontera de decisi\'on entre dos s\'imbolos adyacentes $s_i$ y $s_j$ se define como:

\[
\boxed{\text{La mediatriz perpendicular del segmento que une los dos s\'imbolos}}
\]

\textbf{Formulaci\'on matem\'atica:}

La frontera es el conjunto de puntos $z$ que equidistan de ambos s\'imbolos:
\[
\{z \in \mathbb{C} : |z - s_i| = |z - s_j|\}
\]

Equivalentemente:
\[
\{z \in \mathbb{C} : |z - s_i|^2 = |z - s_j|^2\}
\]

Desarrollando:
\[
|z|^2 - 2\text{Re}(z s_i^*) + |s_i|^2 = |z|^2 - 2\text{Re}(z s_j^*) + |s_j|^2
\]

Simplificando:
\[
\text{Re}[z(s_i^* - s_j^*)] = \frac{|s_i|^2 - |s_j|^2}{2}
\]

Esta es la ecuaci\'on de una recta perpendicular al segmento $[s_i, s_j]$ que pasa por su punto medio si $|s_i| = |s_j|$.

\textbf{Justificaci\'on:}

En un canal AWGN, la decisi\'on ML equivale a elegir el s\'imbolo m\'as cercano en distancia euclidiana al punto recibido. La mediatriz es precisamente el lugar geom\'etrico donde la distancia a ambos s\'imbolos es igual.

\textbf{b-ii) Regiones de decisi\'on para 16-APSK:}

Para 16-APSK (Amplitude and Phase Shift Keying), la constelaci\'on t\'ipicamente tiene dos anillos:
\begin{itemize}
    \item Anillo interior: 4 s\'imbolos con radio $r_1$
    \item Anillo exterior: 12 s\'imbolos con radio $r_2 > r_1$
\end{itemize}

\textbf{Construcci\'on de las regiones:}

\begin{enumerate}
    \item \textbf{Entre s\'imbolos del mismo anillo:}
    
    La frontera es una recta radial (pasa por el origen) que biseca el \'angulo entre los dos s\'imbolos.
    
    Ejemplo: entre s\'imbolos en $0°$ y $90°$ del anillo interior, la frontera es la recta en $45°$.
    
    \item \textbf{Entre s\'imbolos de anillos diferentes:}
    
    La frontera es una recta perpendicular al segmento que une ambos s\'imbolos, pasando por su punto medio.
    
    Esta frontera generalmente no pasa por el origen y tiene forma m\'as compleja.
\end{enumerate}

\textbf{Resultado:}

Cada s\'imbolo queda rodeado por una regi\'on poligonal (generalmente pentagonal o hexagonal para s\'imbolos interiores, y en forma de ''cu\~na'' para s\'imbolos exteriores).

\textbf{Ilustraci\'on conceptual:}

\textit{[En el diagrama, se dibujar\'ian:]}
\begin{itemize}
    \item 16 puntos distribuidos en dos c\'irculos conc\'entricos
    \item Rectas radiales entre s\'imbolos del mismo anillo
    \item Rectas no radiales (mediatrices) entre s\'imbolos de anillos diferentes
    \item Las regiones resultantes forman teselas de Voronoi que cubren todo el plano complejo
\end{itemize}

\textbf{Propiedad importante:}

Las regiones de decisi\'on forman una \textbf{teselaci\'on de Voronoi} del plano complejo, donde cada regi\'on contiene todos los puntos m\'as cercanos a un s\'imbolo particular que a cualquier otro.



\vspace{1cm}
\begin{center}
\rule{0.8\linewidth}{0.4pt}\\
\textit{Fin del examen}
\end{center}

\end{document}
